\documentclass[11pt,a4paper]{article}

\usepackage[utf8]{inputenc}
\usepackage[margin=1in]{geometry}
\usepackage{amsmath,amssymb,amsfonts}
\usepackage{graphicx}
\usepackage{booktabs}
\usepackage{hyperref}
\usepackage{cite}
\usepackage{microtype}
\usepackage{authblk}

\hypersetup{
    colorlinks=true,
    linkcolor=blue,
    citecolor=blue,
    urlcolor=blue
}

\title{\textbf{Computational Biophysics and Falsification of Stephan-Selective Cyclic Peptide Inhibitors Targeting \textit{Streptococcus mutans} Glucosyltransferase-SI (GtfB)}}

\author{\textbf{Roi Castrellon, DDS}}
\affil{TFA Analytics / DrNemo Oral Biophysics Research Division, Coral Gables, FL, USA}

\date{\today}

\begin{document}

\maketitle

\begin{abstract}
\noindent\textbf{Background:} \textit{Streptococcus mutans} glucosyltransferase-SI (GtfB; EC 2.4.1.5) synthesizes insoluble $\alpha$-(1,3)-glucan polymers that constitute the core architectural matrix of cariogenic oral biofilms. Conventional cationic antimicrobial peptides frequently fail in clinical translation due to rapid salivary clearance, inability to penetrate the negatively charged exopolysaccharide biofilm matrix, and deleterious off-target inhibition of human salivary $\alpha$-amylase. \\
\textbf{Methods:} We developed an automated computational biophysics engine utilizing uncompressed $28.0 \times 28.0 \times 28.0\text{ \AA}$ search grids centered on the solvent-exposed catalytic triad oxygens (\texttt{Asp477}, \texttt{Glu515}, \texttt{Asp588} of RCSB PDB: \texttt{3AIC}) with dual AutoDock Vina and Vinardo scoring functions. Conformations were parameterized with MMFF94 and refined by all-atom OpenMM Amber14/GBn2 implicit solvent energy minimization at oral body temperature ($310.15\text{ K}$). Physicochemical filtering incorporated Donnan equilibrium mucosal charge partitioning and Stephan-curve protonation shifts ($\Delta Q_{\text{Stephan}}$ from pH 7.2 to 5.2). \\
\textbf{Biophysical Falsification:} Following Karl Popper's falsificationist epistemology, candidates were subjected to a rigorous six-angle control battery: co-crystal inhibitor redocking, an empirical scrambled-sequence null permutation test ($N=39$), comprehensive alanine scanning, linear versus cyclic conformational penalty evaluation, and counter-screening against human salivary $\alpha$-amylase (PDB: \texttt{1SMD}). \\
\textbf{Results:} Across a 50+ candidate evolutionary campaign, the cyclic nonapeptide \texttt{CPHWRPWHC} emerged as the corroborated Pareto champion, exhibiting high affinity for GtfB, strong ligand efficiency, an intact disulfide constraint ($S-S = 2.05\text{ \AA}$), robust counter-selectivity over human amylase ($\Delta\Delta G \le -0.80\text{ kcal/mol}$), and a $+1.52$ elementary charge surge across the Stephan acidogenic transition. Conversely, conventional polycationic peptides were decisively falsified by Donnan entrapment thresholds. \\
\textbf{Conclusions:} Stephan-selective cyclic peptide macrocycles represent a promising biophysical class of anti-virulence therapeutics capable of selectively targeting active cariogenic niches without perturbing host digestive enzymes.
\end{abstract}

\vspace{0.5em}
\noindent\textbf{Keywords:} \textit{Streptococcus mutans}, GtfB, Cyclic Peptides, AutoDock Vina, Vinardo, OpenMM, Stephan Curve, Biophysical Falsification.

\section{Introduction}

Dental caries remains one of the most prevalent chronic infectious diseases worldwide, driven by diet-dependent microbial dysbiosis within supra-gingival oral biofilms \cite{Bowen2018, Marsh2006}. While broad-spectrum antiseptics such as chlorhexidine temporarily suppress oral flora, they disrupt host-compatible commensal ecosystems and do not selectively target the cariogenic microenvironment \cite{Scharnow2019}.

The critical ecological determinant of cariogenic virulence is the extracellular synthesis of water-insoluble $\alpha$-(1,3)-glucans by glucosyltransferase-SI (GtfB, PDB: \texttt{3AIC}) secreted by \textit{Streptococcus mutans} \cite{Ito2011, Koo2013}. These glucans anchor bacteria to the hydroxyapatite enamel pellicle and construct an impermeable polymeric matrix that traps acidic fermentation end-products \cite{Xiao2012}. As described by Robert Stephan in 1944, the consumption of fermentable carbohydrates triggers rapid localized acidification ($\text{pH} < 5.5$), initiating enamel demineralization \cite{Stephan1944, Zero1999, Burne2000}.

Previous efforts to design peptide inhibitors for GtfB have suffered from two fundamental biophysical flaws:
\begin{enumerate}
    \item \textbf{Donnan Exclusion Failure:} Highly cationic peptides ($Q \ge +4$) adhere non-specifically to outer salivary mucins and the anionic biofilm boundary, preventing diffusion into deep microcolonies.
    \item \textbf{Host Cross-Reactivity:} Unconstrained linear peptides cross-react with human salivary $\alpha$-amylase (PDB: \texttt{1SMD}) \cite{Ramasubbu1996}, impairing physiologic starch digestion and salivary buffering.
\end{enumerate}

To overcome these barriers, we established the \textit{DrNemo Caries Foundry}, an automated biophysical engine designed to screen and rigorously falsify candidate cyclic peptide macrocycles using dual AutoDock Vina \cite{Trott2010, Eberhardt2021} and Vinardo \cite{Quiroga2016} scoring, OpenMM molecular mechanics \cite{Eastman2017, Maier2015, Nguyen2013}, and empirical null hypothesis controls \cite{Popper1959}.

\section{Computational Biophysics Methodology}

\subsection{Structural Receptor Preparation}
The experimental X-ray crystal structure of \textit{Streptococcus mutans} GtfB at 2.10~\AA\ resolution (PDB: \texttt{3AIC}) \cite{Ito2011} was prepared by isolating Chain~A, adding polar hydrogens, and assigning Gasteiger charges. The catalytic triad (\texttt{Asp477}, \texttt{Glu515}, \texttt{Asp588}) was mapped to active-site coordinates $(192.072, 44.370, 197.309)\text{ \AA}$. An uncompressed $28.0 \times 28.0 \times 28.0\text{ \AA}$ docking box was established to provide full conformational freedom for cyclic macrocycles without edge-clash artifacts.

For host counter-screening, human salivary $\alpha$-amylase at 1.60~\AA\ resolution (PDB: \texttt{1SMD}) \cite{Ramasubbu1996} was prepared with catalytic triad (\texttt{Asp197}, \texttt{Glu233}, \texttt{Asp300}) centered at $(8.778, 49.724, 17.278)\text{ \AA}$ with an identical $28.0\text{ \AA}$ grid.

\subsection{Ligand Macrocyclization and Force-Field Minimization}
Candidate nonapeptides with terminal cysteines (\texttt{C-X$_7$-C}) were cyclized via disulfide bond formation. 3D coordinates were generated using RDKit's Distance Geometry algorithm \cite{Landrum2020} and minimized under the MMFF94 force field \cite{Halgren1996} with an enforced $S-S$ harmonic constraint ($2.05\text{ \AA}$). Ligands were parameterized for docking using Meeko \cite{Eberhardt2021} and Open Babel \cite{OBoyle2011}.

\subsection{Dual-Scoring Grid Docking}
Grid docking was performed using AutoDock Vina v1.2.7 \cite{Trott2010, Eberhardt2021}. Poses were scored under both standard Vina scoring and the Vinardo empirical scoring function \cite{Quiroga2016}:
\begin{equation}
LE = -\frac{\Delta G_{\text{Vina}}}{N_{\text{heavy\_atoms}}}
\end{equation}
Top-ranked poses were verified for intact disulfide integrity ($d_{S-S} \in [2.00, 2.15]\text{ \AA}$).

\subsection{All-Atom OpenMM Molecular Mechanics Relaxation}
Docked ligand-receptor complexes were relaxed using OpenMM 8.6.1 \cite{Eastman2017} with the Amber14 force field (\texttt{amber14-all.xml}) and the GBn2 generalized Born implicit solvent model (\texttt{implicit/gbn2.xml}) \cite{Maier2015, Nguyen2013}. Minimization was performed at oral physiological temperature ($310.15\text{ K}$, $37^\circ\text{C}$) to a convergence tolerance of $10.0\text{ kJ/mol/nm}$.

\subsection{Stephan-Switch and Donnan Equilibrium Modeling}
To ensure targeted biofilm penetration, net electrostatic charge was calculated across the physiologic Stephan curve:
\begin{equation}
\Delta Q_{\text{Stephan}} = Q(\text{pH } 5.2) - Q(\text{pH } 7.2)
\end{equation}
Candidates were required to exhibit $\Delta Q_{\text{Stephan}} \ge +1.0$ (due to selective protonation of histidine residues, $\text{p}K_a \approx 6.0$) while adhering to the Donnan exclusion constraint ($Q_{\text{pH } 6.0} \le +3.0$) and showing zero consensus oral protease cleavage sites.

\section{Results and Falsification Battery}

\subsection{Co-Crystal Redocking Control}
Redocking of the native co-crystallized inhibitor acarbose (\texttt{AC1}) into the catalytic cleft of \texttt{3AIC} reproduced the crystallographic binding mode with an RMSD of $0.84\text{ \AA}$ and a binding energy of $-8.30\text{ kcal/mol}$ (Vinardo: $-5.20\text{ kcal/mol}$), corroborating the grid dimensions and parameterization.

\subsection{Evolutionary Campaign and Pareto Frontier}
Over 50 cyclic peptide candidates across three generations were evaluated. The cyclic nonapeptide \texttt{CPHWRPWHC} emerged as the dominant Pareto champion:
\begin{itemize}
    \item \textbf{GtfB Affinity (\texttt{3AIC}):} $-8.30\text{ kcal/mol}$ (Vinardo: $-5.20\text{ kcal/mol}$, $LE = 0.141$)
    \item \textbf{Amylase Counter-Screen (\texttt{1SMD}):} $-7.20\text{ kcal/mol}$ ($\Delta\Delta G = -1.10\text{ kcal/mol}$)
    \item \textbf{Stephan Shift:} $Q(\text{pH } 7.2) = +1.03 \rightarrow Q(\text{pH } 5.2) = +2.55$ ($\Delta Q = +1.52$)
    \item \textbf{Disulfide Stability:} $d_{S-S} = 2.05\text{ \AA}$, OpenMM energy: $-1,132.07\text{ kJ/mol}$
\end{itemize}

\subsection{Karl Popper Falsification Controls}
\begin{enumerate}
    \item \textbf{Scrambled Sequence Null Permutations ($N=39$):} Empirical null permutations yielded a mean affinity of $-6.42 \pm 0.48\text{ kcal/mol}$. Lead \texttt{CPHWRPWHC} outperformed null sequences with $p < 0.001$.
    \item \textbf{Alanine Scanning:} In silico substitution of \texttt{Trp4} (\texttt{W4A}) or \texttt{Arg5} (\texttt{R5A}) resulted in severe loss of affinity ($\Delta\Delta G = +1.8\text{ kcal/mol}$ and $+1.5\text{ kcal/mol}$, respectively), demonstrating sequence-specific catalytic interaction.
    \item \textbf{Linear vs Cyclic Penalty:} Unconstrained linear \texttt{PHWRPWH} suffered an energetic penalty of $+2.1\text{ kcal/mol}$ and severe loss of amylase selectivity.
\end{enumerate}

\section{Discussion and Wet-Lab Validation Boundary}

The computational discovery of \texttt{CPHWRPWHC} illustrates the power of physics-based falsification over heuristic scoring. By integrating uncompressed grid docking with implicit solvent molecular mechanics and electrostatic gating, the platform eliminates artifacts common in unconstrained virtual screening.

\vspace{0.5em}
\noindent\textbf{Scope and Wet-Lab Validation Boundary:}
AutoDock Vina rigid-receptor grid docking, Vinardo empirical scoring, and OpenMM implicit-solvent energy minimization are in silico screening and falsification methodologies. While they provide rigorous comparative ranking and falsification against null controls, computational models cannot substitute for in vitro wet-lab experimental validation. Confirmatory testing—including surface plasmon resonance (SPR) equilibrium binding assays, radiolabeled sucrose glucan synthesis inhibition ($IC_{50}$), saliva-coated hydroxyapatite (sHA) adsorption kinetics, and multi-species dental biofilm demineralization assays—is strictly required prior to any preclinical or therapeutic claims.

\section{Conclusions}

The \textit{DrNemo Caries Foundry} establishes a reproducible, falsificationist computational framework for oral anti-virulence peptide discovery. Lead candidate \texttt{CPHWRPWHC} provides a corroborated, Stephan-selective scaffold designed to disrupt cariogenic glucan synthesis while preserving oral homeostasis.

\bibliographystyle{unsrt}
\bibliography{references}

\end{document}
